\documentclass[oneside,12pt]{wipb}

\katedra{data science}
\typpracy{magisterska}
\temat{Multi-Agent SRE: Uczenie w Ewoluującym Kontekście}
\subject{Multi-Agent SRE: Learning in Evolving Context}
\autor{Michał Murawski}
\promotor{dr inż. Mateusz Żbikowski}
\supervisor{dr inż. Mateusz Żbikowski}
\indeks{s34597}
\studia{niestacjonarne}
\profil{studia II stopnia}
\kierunekstudiow{informatyka}
\specjalnosc{data science}
\specialisation{data science}

\hypersetup{
pdfauthor={Imię Nazwisko},
pdftitle={Praca inżynierska},
pdfsubject={Temat pracy},
pdfkeywords={programowanie, bluetooth, jazda konna},
pdfpagemode=UseNone,
linkcolor=black,
citecolor=black,
urlcolor=black
} 

\setlength{\epigraphwidth}{1\textwidth}

\begin{document}

\maketitle
\pagestyle{plain}

\chapter*{Streszczenie}

Celem niniejszej pracy magisterskiej jest eksperymentalna ocena wpływu mechanizmu adaptacyjnej ewolucji kontekstu agentowego na skuteczność autoremediacji aplikacji mikrousługowych w środowisku Kubernetes. W pracy zaprojektowano hierarchiczny system wieloagentowy oparty na paradygmacie \textit{agent-as-a-tool}, w którym \textit{Incident Commander} koordynuje pracę trzech wyspecjalizowanych deputatów, wyposażony w mechanizm samouczenia oparty na algorytmie \textit{Agentic Context Engineering} (ACE), rozszerzonym o hierarchiczną propagację refleksji wzdłuż drzewa zadań oraz mechanizm ukrywania wag skuteczności reguł przed agentem wykonawczym.

Eksperymenty przeprowadzono na izolowanym środowisku z referencyjną aplikacją mikrousługową \textit{Robot-shop} i autorskim komponentem \textit{Saboteur Agent}, generującym deterministyczną sekwencję 98 incydentów reprezentujących cztery klasy usterek. Badania objęły dwa modele językowe - \textit{GPT-5-nano} i \textit{Minimax-2.5} - przy identycznej sekwencji epizodów, z automatyczną oceną trzech binarnych wskaźników sukcesu.

Wyniki potwierdzają główną tezę pracy: \textit{GPT-5-nano} osiąga poprawę skuteczności remediacji o $+18{,}4$~pp ($+120\%$ względnie), \textit{Minimax-2.5} o $+5{,}1$~pp ($+13\%$ względnie). Modele wykształcają odmienne strategie organizacyjne - GPT-5-nano uczy się koordynacji agentów, Minimax-2.5 uczy się zamykania cyklu remediacji przez wdrożenie. Zaobserwowano zjawisko emergentne: średnia maksymalna głębokość drzewa zadań rośnie o $+26{,}1\%$ bez jawnego programowania tej strategii.

\vspace{20pt}
\noindent Słowa kluczowe: \textit{Site Reliability Engineering}, autoremediacja, systemy wieloagentowe, duże modele językowe, \textit{Agentic Context Engineering}, samouczenie agentów, mikrousługi, Kubernetes, drzewo zadań


\setcounter{tocdepth}{2}
\tableofcontents

\chapter{Wstęp}
\section{Kontekst i tło problemu}
Współczesne systemy informatyczne opierają się w przeważającej mierze na architekturze rozproszonych mikrousług, których orkiestracja odbywa się przy pomocy platform takich jak Kubernetes. Skala tych środowisk nieustannie rośnie – dobitnym przykładem są globalne przedsiębiorstwa technologiczne, takie jak Netflix, których platformy operują na zawiłej sieci ponad 1000 ściśle współpracujących ze sobą mikrousług.

Choć podejście oparte na konteneryzacji przyniosło ogromne korzyści w zakresie skalowalności, elastyczności i niezależności wdrożeń, drastycznie zwiększyło ono również złożoność operacyjną środowisk IT. Zapewnienie ciągłej operacyjności (High Availability) tak skomplikowanych systemów rozproszonych pochłania gigantyczne nakłady finansowe. Koszty te wynikają w głównej mierze z konieczności opłacania i utrzymywania wysoce wykwalifikowanych zespołów inżynierskich (Site Reliability Engineering / DevOps), które zmuszone są do pracy w systemie całodobowych dyżurów (tzw. "on-call").

Kiedy dochodzi do awarii, presja czasu jest ogromna. Ręczna diagnoza i przywrócenie systemu do stanu używalności (remediacja) wymagają od dyżurujących inżynierów analizy setek powiązań sieciowych, rozproszonych logów, metryk i konfiguracji klastra. 
Stanowi to nie tylko rosnące obciążenie dla zespołów IT, ale również realne straty biznesowe z powodu przedłużających się przestojów.

Kluczowym rozwiązaniem tego narastającego problemu staje się zaawansowana automatyzacja procesów naprawczych. Jak zauważają autorzy pracy "Leveraging Large Language Models for the Auto-remediation of Microservice Applications: An Experimental Study", Duże Modele Językowe (LLM) wykazują ogromny potencjał w automatyzacji zadań związanych z diagnozowaniem i rozwiązywaniem incydentów w mikrousługach. 
Jednakże standardowe, jednoagentowe podejścia często napotykają na bariery przy rosnącej złożoność systemu informatycznego i brakiem specjalistycznej wiedzy o konkretnym klastrze czy niemożnością utrzymania długoterminowego kontekstu zdarzeń. 

\section{Cel badawczy i wykorzystanie systemów wieloagentowych}
Aby przezwyciężyć ograniczenia pojedynczych modeli, coraz częściej sięga się po systemy wieloagentowe (Multi-Agent Systems – MAS). 
Wdrożenie takich systemów wiąże się jednak z licznymi wyzwaniami dotyczącymi koordynacji, komunikacji i stabilności działania, co szeroko omówiono w publikacji 
"LLM Multi-Agent Systems: Challenges and Open Problems". 
W krytycznych zastosowaniach, takich jak zarządzanie infrastrukturą, kluczowe jest zapewnienie rzetelności podejmowanych działań, co podkreślono w badaniach nad 
"Reliable Decision-Making for Multi-Agent LLM System".

Głównym celem niniejszej pracy jest eksperymentalne zbadanie, w jaki sposób system wieloagentowy może zwiększyć skuteczność i niezawodność auto-remediacji aplikacji w środowisku Kubernetes. 
Aby zapewnić ciągłe uczenie się systemu i minimalizację błędów, w pracy zaimplementowano nowatorskie podejście Inżynierii Kontekstu – potok Agentic Context Evolution (ACE), 
inspirowany koncepcją przedstawioną w badaniu "Agentic Context Engineering: Evolving Contexts for Self-Improving Language Models". Metoda ta pozwala agentom na dynamiczne tworzenie i 
modyfikowanie własnej bazy wiedzy (Playbook/SOP) w miarę rozwiązywania kolejnych incydentów.

\section{Teza i pytania badawcze}
Główna teza pracy zakłada, że zastosowanie wieloagentowego systemu opartego na potoku Agentic Context Evolution (ACE) 
znacząco poprawia skuteczność procesów auto-remediacji w klastrach Kubernetes, skracając czas naprawy (TTR - z ang. Time To Resolution) i zwiększając ogólny wskaźnik rozwiązywalności 
(Success Rate) w porównaniu do systemów o statycznym kontekście.


W celu weryfikacji postawionej tezy, sformułowano następujące pytania badawcze: \newline

1. W jaki sposób dynamiczna ewolucja instrukcji operacyjnych (Playbooka) wpływa na skuteczność agentów w obliczu różnorodnych klas incydentów (od problemów infrastrukturalnych po błędy w kodzie aplikacji)?

2. Jak zmiana częstotliwości aktualizacji kontekstu (przejście od fazy eksploracji do eksploatacji, analogicznie do strategii $\epsilon$-greedy z uczenia ze wzmocnieniem) wpływa na stabilność i jakość generowanych heurystyk?

3. Czy implementacja mechanizmu "Reflektora" w pętli ACE skutecznie zapobiega regresjom wprowadzanym przez nowe, błędne reguły naprawcze?

\section{Zakres i struktura pracy}

Praca skupia się na stworzeniu autorskiego środowiska eksperymentalnego, symulującego rzeczywiste awarie w klastrze Kubernetes. Incydenty zostały podzielone na pięć klas (m.in. awarie sieciowe wywoływane przez Chaos Mesh, błędy w kodzie, czy złożone scenariusze typu "Red Herring").

W kolejnych rozdziałach praca przedstawia: przegląd aktualnego stanu wiedzy w obszarze LLM i operacji IT (Rozdział 2),
 szczegółowy opis architektury potoku ACE i metodologii badawczej (Rozdział 3), projekt przeprowadzonych eksperymentów wraz ze scenariuszami testowymi (Rozdział 4),
  a także rzetelną ewaluację zebranych metryk wydajnościowych, 
takich jak Time To Resolution czy jakość wypracowanych heurystyk (Rozdział 5). Pracę zamyka podsumowanie wyników oraz wskazanie kierunków dalszych badań.

Piękno róży przyćmiło moją jaźń \cite{cleancode, booking:_nodate}. 


%\nocite{*}
\bibliographystyle{unsrt}
\bibliography{bibliografia}
\addcontentsline{toc}{chapter}{Bibliografia}

%\listoffigures
{
    \let\oldnumberline\numberline%
    \renewcommand{\numberline}{\figurename~\oldnumberline}%
    \listoffigures
}
\addcontentsline{toc}{chapter}{Lista figur}

\end{document}


